% Ad Atticum 4.9 (ad-atticum-04-09)
% Source: https://raw.githubusercontent.com/PerseusDL/canonical-latinLit/master/data/phi0474/phi057/phi0474.phi057.perseus-lat1.xml
% Fetched by scripts/fetch_latin.py

% Dateline (Perseus): Scr. Neapoli iv K. Mai a. 699 (55).

\ciceroLetterOpener{CICERO ATTICO salutem}

\ciceroSection{1} sane velim scire num censum impediant tribuni diebus vitiandis (est enim hic rumor) totaque de censura quid agant, quid cogitent. nos hic cum Pompeio fuimus. multa mecum de re publica sane sibi displicens, ut loquebatur (sic est enim in hoc homine dicendum), Syriam spernens, Hispaniam iactans, hic quoque, ut loquebatur; et opinor usque quaque, de hoc cum dicemus, sit hoc quasi καὶ τόδε . tibi etiam gratias agebat quod signa componenda suscepisses; in nos vero suavissime hercule est effusus. venit etiam ad me in Cumanum a se. nihil minus velle mihi visus est quam Messallam consulatum petere. de quo ipso si quid scis velim scire.

\ciceroSection{2} quod Lucceio scribis te nostram gloriam commendaturum et aedificium nostrum quod crebro invisis, gratum. Quintus frater ad me scripsit se, quoniam Ciceronem suavissimum tecum haberes, ad te Nonis Maias venturum. ego me de Cumano movi ante diem v Kal. Maias. eo die Neapoli apud Paetum. ante diem IIII Kal. Maias iens in Pompeianum bene mane haec scripsi.
